\section{Verification and scope}\label{sec:verification}

The all-length multiplier theorem and transversal-group dichotomy are proved
in the text and use no finite enumeration.  The only imported ingredients in
the group classification are the stabilizer-AME rigidity and
stabilizer-character correction stated in the introduction
\cite{RuddAMEStabilizerRigidity2026}.

\begin{table}[p]
\centering
\small
\begin{tabular}{@{}L{.20\textwidth}L{.17\textwidth}L{.27\textwidth}L{.25\textwidth}@{}}
\toprule
Claim family & Status & Verification method & Imported or trusted input \\
\midrule
Multiplier nullity & proved here & MDS distance and a vanishing-coordinate argument & none \\
Exact fixed-coordinate group & proved here & block preservation, explicit unipotents, and logical Paulis & imported product-unitary rigidity and phase correction \\
Six-point conic criterion & proved here & MDS--CSS dictionary and Gale-transform argument & none \\
Pencil quotient & proved here & symbolic bracket identities and exact transition-map enumeration & imported minimum-support transition theorem; certificate replay for the enumeration \\
Frobenius-sector formulas & proved here & direct finite-field algebra & no complete extension-field orbit classification \\
Clebsch application & consequence & syndrome lemma composed with the cited Clebsch deep-hole theorem & cited conic, count, and orbit result \\
Scalar separators and transport divisor & appendix results & exact generators, JSON certificates, and independent replay & finite exact computation \\
Coordinate-permutation extensions & twelve computed cases & exhaustive permutation checks and explicit complements & no general classification beyond the listed cases \\
Formalization & partial & Lean kernel checking through the paper-specific gate & hypothesis-explicit interfaces and three native finite cardinality evaluations \\
\bottomrule
\end{tabular}
\caption{Proof and verification status of the main claim families.  A
hypothesis-explicit interface proves a conclusion from named mathematical or
computational inputs; importing it does not prove those inputs.}
\label{tab:verification-status}
\end{table}

The paper-local supplement contains the exact generators, compact JSON
certificates, and a deterministic standard-library verifier for the six-point
computations.  Running
\begin{center}
\texttt{python3 supplement/verify.py --replay}
\end{center}
checks the manifest and regenerates every certificate in memory.  The
supplement records each finite domain, stopping condition, independent replay,
and negative scope statement.

The Lean~4 development uses toolchain \texttt{v4.32.0-rc1} with a pinned
Mathlib revision.  Its paper-specific roots are
\begin{center}
\texttt{RelativeConicArcs.Gates.MDSCSSTransversalGeometry}\
\texttt{RelativeConicArcs.Gates.MDSCSSTransversalGeometryAxioms}.
\end{center}
The first is the import-only gate for the checked surface; the second prints
the axiom dependencies of every cited declaration.  Kernel reduction covers
the MDS--CSS dictionary, multiplier nullity, pencil quotient,
extension-field scalar formulas, syndrome geometry, contraction algebra,
transport cycle algebra, and the abstract consequences of a supplied
coordinate-permutation complement.  The exact group, finite classifications,
separators, and transport geometry use interfaces whose input fields are
proved in the text or supplied by the certificates.  Three graph
cardinalities use exhaustive native evaluation; two marginal-moment results
inherit the star-count evaluation axiom.  Apart from those local declarations,
the cited formal results use only propositional extensionality, choice, and
quotient soundness.

The local-unitary classification by \(z\) is restricted to odd prime fields.
Over extension fields, Frobenius is already an additional local Clifford and
the full additive-Clifford orbit problem is not claimed.  The exact group
theorem fixes the physical coordinates and assumes a product implementation.
It does not classify arbitrary coordinate-permutation enlargements,
non-product physical implementations, nonstabilizer AME states, or unrelated
quantum-code families.
